\chapter{Arithmetic Units}
\label{ch:arith}

We can select (Chapter~\ref{ch:mux}) and distribute signals, but so far we work
only with yes/no bits. A question remains: how does a circuit made of gates
\emph{calculate}? The word ``computer'' means ``one that computes'', and the
main thing a computer does is calculate. In this chapter we will see that
adding, subtracting, and comparing numbers are done by the very same logic
gates, and then we will assemble them into an arithmetic logic unit.

\begin{buildbox}
\textbf{Task.} Add two one-bit numbers and get a sum and a carry.

\textbf{In Logra.} Learn tab, the ``Circuits and the board'' book, the
``Arithmetic'' chapter, the lesson ``A half adder from gates'' (Designer): you
need two gates, XOR and AND.

\textbf{Expected result.} For inputs $a,b$ the sum output equals $a\oplus b$ and
the carry output equals $a\cdot b$; for $1+1$ you get sum 0 and carry 1. The
rest of the chapter shows how to extend this circuit to multibit numbers,
negative numbers, and comparison.
\end{buildbox}

\section{Adding Binary Numbers}

Binary numbers are added in columns, like decimal ones, except that a carry to
the next position occurs when the sum reaches two:
\[
  0+0 = 0, \quad 0+1 = 1, \quad 1+0 = 1, \quad 1+1 = 10_2 \ (\text{zero and a carry of one}).
\]

\begin{figure}[H]
\centering
\begin{tikzpicture}[x=9mm, y=7mm, font=\ttfamily\large]
  \node[text=FpgaRed, font=\ttfamily\small] at (1,3) {1}; \node[text=FpgaRed, font=\ttfamily\small] at (2,3) {1};
  \node[text=FpgaRed, font=\ttfamily\small] at (3,3) {1};
  \foreach \d [count=\i] in {0,1,1,1} \node at (\i,2) {\d};
  \foreach \d [count=\i] in {0,0,1,1} \node at (\i,1) {\d};
  \foreach \d [count=\i from 0] in {1,0,1,0} \node at (\i+1,-0.2) {\d};
  \node at (-0.2,1.5) {$+$};
  \draw[thick] (0.4,0.5) -- (4.5,0.5);
  \node[font=\sffamily\small, anchor=west] at (5.2,2) {$0111_2 = 7$};
  \node[font=\sffamily\small, anchor=west] at (5.2,1) {$0011_2 = 3$};
  \node[font=\sffamily\small, anchor=west] at (5.2,-0.2) {$1010_2 = 10$};
  \node[font=\sffamily\small, anchor=west, text=FpgaRed] at (5.2,3) {carries};
\end{tikzpicture}
\caption{Column addition in the binary system}
\end{figure}

\begin{diybox}[Column addition]
\textbf{Task.} \stars{1} Add $0111_2+0011_2$ on paper and confirm the carry
between bit positions. \stars{2} Predict for which $a,b$ the top carry occurs in
the $4$-bit sum $a+b$, and check in Logra.

\textbf{In Logra.} Learn tab, the ``Arithmetic'' chapter, the lesson ``An adder
as one block''; then Waveforms.

\textbf{Expected result.} The sum $7+3=10=1010_2$ with no carry out of the top
bit; $15+1$ gives zero and carry 1.

\textbf{How to check.} Waveforms: apply $7$ and $3$, then $15$ and $1$, and
compare the sum and carry-output values with your paper calculation.
\end{diybox}

\section{The Half Adder}

Adding two one-bit numbers $a$ and $b$ gives a two-bit result: a sum bit $s$ and
a carry bit $c$. The truth table:

\begin{figure}[H]
\centering
\begin{minipage}{0.35\textwidth}
\centering
\begin{tabular}{cc|cc}
\toprule
$a$ & $b$ & $c$ & $s$ \\
\midrule
0 & 0 & 0 & 0 \\
0 & 1 & 0 & 1 \\
1 & 0 & 0 & 1 \\
1 & 1 & 1 & 0 \\
\bottomrule
\end{tabular}

\medskip
$s = a \oplus b, \qquad c = a\cdot b$
\end{minipage}
\begin{minipage}{0.55\textwidth}
\centering
\begin{tikzpicture}
  \node[gXOR, minimum height=12mm] (x) at (2,0.8) {};
  \node[gAND, minimum height=12mm] (g) at (2,-0.8) {};
  \draw[wire] (x.input 1) -- ++(-2.0,0) node[left] {$a$};
  \draw[wire] (x.input 2) -- ++(-2.0,0) node[left] {$b$};
  \coordinate (A) at ($(x.input 1)+(-1.4,0)$); \coordinate (B) at ($(x.input 2)+(-0.9,0)$);
  \draw[wire] (A) |- (g.input 1); \fill (A) circle (1.3pt);
  \draw[wire] (B) |- (g.input 2); \fill (B) circle (1.3pt);
  \draw[wire] (x.output) -- ++(0.6,0) node[right] {$s$ (sum)};
  \draw[wire] (g.output) -- ++(0.6,0) node[right] {$c$ (carry)};
\end{tikzpicture}
\end{minipage}
\caption{The half adder}
\label{fig:ha}
\end{figure}

\begin{diybox}[The half adder]
\textbf{Task.} \stars{1} Build a half adder from an XOR gate and an AND gate.
\stars{2} Connect the outputs to two lamps and check all four combinations.

\textbf{In Logra.} Learn tab, the ``Arithmetic'' chapter, the lesson ``A half
adder from gates''; the reference is the example \code{half-adder}.

\textbf{Expected result.} The sum $s$ and carry $c$ as in the table above.

\textbf{How to check.} Waveforms: the four combinations of $a,b$; RTL: two gates
on the inputs $a,b$, as in Fig.~\ref{fig:ha}.
\end{diybox}

\section{The Full Adder}

A half adder is good only for the lowest bit. In the other bit positions
\emph{three} bits are added: $a_i$, $b_i$, and the carry $c_\text{in}$ from the
previous position. Such a unit is called a \term{full adder}.

\begin{figure}[H]
\centering
\begin{minipage}{0.38\textwidth}
\centering
\small
\begin{tabular}{ccc|cc}
\toprule
$a$ & $b$ & $c_\text{in}$ & $c_\text{out}$ & $s$ \\
\midrule
0&0&0 & 0&0 \\
0&0&1 & 0&1 \\
0&1&0 & 0&1 \\
0&1&1 & 1&0 \\
1&0&0 & 0&1 \\
1&0&1 & 1&0 \\
1&1&0 & 1&0 \\
1&1&1 & 1&1 \\
\bottomrule
\end{tabular}
\end{minipage}
\begin{minipage}{0.58\textwidth}
\[
  s = a \oplus b \oplus c_\text{in}
\]
\[
  c_\text{out} = ab + ac_\text{in} + bc_\text{in}
\]
\small Note that $c_\text{out}$ is the ``2 out of 3'' majority function from
Chapter~\ref{ch:comb}! A carry occurs when there are at least two ones.
\end{minipage}
\caption{The full adder: truth table and formulas}
\end{figure}

\begin{figure}[H]
\centering
\begin{tikzpicture}[ha/.style={draw=FpgaBlue, fill=FpgaLight, thick, minimum width=14mm, minimum height=14mm, font=\sffamily\small, align=center}]
  \node[ha] (h1) at (0,0) {HA};
  \node[ha] (h2) at (2.8,-0.6) {HA};
  \node[gOR, minimum width=10mm, minimum height=7mm] (o) at (5.4,-1.8) {};
  \draw[wire] ([yshift=3.5mm]h1.west) -- ++(-0.8,0) node[left] {$a$};
  \draw[wire] ([yshift=-3.5mm]h1.west) -- ++(-0.8,0) node[left] {$b$};
  \draw[wire] ([yshift=3.5mm]h1.east) -- node[above, lbl] {$s$} ++(0.4,0) |- ([yshift=3.5mm]h2.west);
  \draw[wire] (-0.8,-1.4) node[left] {$c_\text{in}$} -- (1.7,-1.4) |- ([yshift=-3.5mm]h2.west);
  \draw[wire] ([yshift=3.5mm]h2.east) -- ++(3.4,0) node[right] {$s$};
  \draw[wire] ([yshift=-3.5mm]h1.east) -- node[below, lbl, pos=0.3] {$c$} ++(0.25,0) |- (0.95,-2.2) -- (4.7,-2.2) |- (o.input 2);
  \draw[wire] ([yshift=-3.5mm]h2.east) -- node[above, lbl] {$c$} ++(0.4,0) |- (o.input 1);
  \draw[wire] (o.output) -- ++(0.6,0) node[right] {$c_\text{out}$};
\end{tikzpicture}
\caption{A full adder made of two half adders and an OR gate}
\label{fig:fa}
\end{figure}

\begin{diybox}[The full adder]
\textbf{Task.} \stars{1} Build a full adder from two half adders and an OR gate
as in Fig.~\ref{fig:fa}. \stars{2} Build the carry differently: as ``2 out of
3'' ($ab+ac_\text{in}+bc_\text{in}$) and compare the two versions.

\textbf{In Logra.} Learn tab, the ``Arithmetic'' chapter, the lesson ``A full
adder''; the reference is the example \code{full-adder}.

\textbf{Expected result.} All $8$ combinations give the values of the table
$c_\text{out}$, $s$; both ways of getting the carry agree.

\textbf{How to check.} Waveforms: go through all $8$ rows; RTL: compare the
number of gates in the two versions.
\end{diybox}

\section{The Multibit Adder}

To add $n$-bit numbers, $n$ full adders are connected in a chain: the carry
output of each bit position goes to the carry input of the next. Such an adder
is called a \term{ripple-carry adder}.

\begin{figure}[H]
\centering
\begin{tikzpicture}[fa/.style={draw=FpgaBlue, fill=FpgaLight, thick, minimum width=15mm, minimum height=13mm, font=\sffamily\small}]
  \foreach \i [evaluate=\i as \x using (3-\i)*2.9] in {0,...,3} {
    \node[fa] (f\i) at (\x,0) {FA$_\i$};
    \draw[wire] ([xshift=-3mm]f\i.north) -- ++(0,0.6) node[above] {$a_\i$};
    \draw[wire] ([xshift=3mm]f\i.north) -- ++(0,0.6) node[above] {$b_\i$};
    \draw[wire] (f\i.south) -- ++(0,-0.6) node[below] {$s_\i$};
  }
  \foreach \i [evaluate=\i as \j using int(\i+1)] in {0,1,2}
    \draw[arr, FpgaRed] (f\i.west) -- node[above, lbl] {$c_\j$} (f\j.east);
  \draw[arr] ($(f0.east)+(0.8,0)$) node[right] {$c_0 = 0$} -- (f0.east);
  \draw[arr, FpgaRed] (f3.west) -- ++(-0.8,0) node[left, text=FpgaDark] {$c_4$ (carry)};
\end{tikzpicture}
\caption{A four-bit ripple-carry adder. The carry chain is shown in red}
\label{fig:rca}
\end{figure}

The drawback of such an adder: the top bit cannot compute its result until it
receives the carry from the lower bit, and that bit waits for the carry from an
even lower one. The delay grows in proportion to the number of bits. That is why
fast processors use adders with \term{carry lookahead}, in which the carry is
computed in parallel by a special circuit. \lookahead{in an FPGA, a separate
fast \emph{carry chain} is laid next to each logic block for this purpose
(Chapter~\ref{ch:inside}).}

\begin{figure}[H]
\centering
\begin{tikzpicture}
\begin{axis}[
  width=0.8\textwidth, height=5.4cm, xlabel={Width, bits}, ylabel={Delay, $t_\text{d}$},
  xmin=0, xmax=64, ymin=0, ymax=140, xtick={8,16,32,48,64},
  legend pos=north west, legend style={font=\sffamily\small, draw=none, fill=white, fill opacity=0.8},
  label style={font=\sffamily\small}, tick label style={font=\sffamily\small},
  grid=major, grid style={FpgaGray!20}, axis lines*=left,
]
  \addplot[very thick, FpgaRed, domain=1:64, samples=2] {2*x};
  \addlegendentry{ripple carry ($\sim n$)}
  \addplot[very thick, FpgaGreen, domain=1:64, samples=64] {4+4*ln(x)/ln(2)};
  \addlegendentry{carry lookahead ($\sim \log_2 n$)}
\end{axis}
\end{tikzpicture}
\caption{Adder delay as a function of width (an estimate)}
\end{figure}

\begin{diybox}[A four-bit adder]
\textbf{Task.} \stars{1} Build a 4-bit adder from four full adders as in
Fig.~\ref{fig:rca}. \stars{2} Check $7+9$ and $15+1$ and explain where $c_4$
comes from. \stars{3} Using problem~5 of this chapter, estimate how long the
carry takes to travel along the chain.

\textbf{In Logra.} Designer; the reference is the lesson ``An adder as one
block'' and the example \code{carry-overflow}.

\textbf{Expected result.} $7+9=0000$ with carry $c_4=1$ (the number 16 does not
fit in 4 bits); $15+1=0000$ with $c_4=1$.

\textbf{How to check.} Waveforms: in a chain of four adders the output $s_3$
changes later than $s_0$; compare the ready-made block and your chain with the
same inputs.
\end{diybox}

\section{Negative Numbers and Subtraction}

How do we write a negative number if we have only zeros and ones? We use
\term{two's complement}: the top bit is considered to have a \emph{negative}
weight. For 4 bits the weights of the bit positions are $-8, 4, 2, 1$.

\begin{table}[H]
\centering
\caption{Four-bit numbers in two's complement}
\small
\begin{tabular}{c c @{\hspace{10mm}} c c}
\toprule
\textbf{Code} & \textbf{Number} & \textbf{Code} & \textbf{Number} \\
\midrule
\code{0000} & 0 & \code{1000} & $-8$ \\
\code{0001} & 1 & \code{1001} & $-7$ \\
\code{0010} & 2 & \code{1010} & $-6$ \\
\code{0011} & 3 & \code{1011} & $-5$ \\
\code{0100} & 4 & \code{1100} & $-4$ \\
\code{0101} & 5 & \code{1101} & $-3$ \\
\code{0110} & 6 & \code{1110} & $-2$ \\
\code{0111} & 7 & \code{1111} & $-1$ \\
\bottomrule
\end{tabular}
\end{table}

\begin{figure}[H]
\centering
\begin{tikzpicture}
  \foreach \i in {0,...,15} {
    \pgfmathsetmacro{\ang}{90-\i*22.5}
    \pgfmathtruncatemacro{\v}{\i<8 ? \i : \i-16}
    \fill[\ifnum\i<8 FpgaBlue\else FpgaOrange\fi] (\ang:2.2) circle (0.09);
    \node[font=\sffamily\small] at (\ang:2.65) {\v};
  }
  \draw[thick, FpgaGray] (0,0) circle (2.2);
  \node[font=\ttfamily\scriptsize] at (90:1.75) {0000};
  \node[font=\ttfamily\scriptsize] at (90-7*22.5:1.75) {0111};
  \node[font=\ttfamily\scriptsize] at (90-8*22.5:1.75) {1000};
  \node[font=\ttfamily\scriptsize] at (90-15*22.5:1.75) {1111};
  \draw[-{Stealth}, thick, FpgaGreen] (60:1.2) arc (60:-30:1.2);
  \node[font=\sffamily\small, text=FpgaGreen] at (0,0) {$+1$};
  \node[font=\sffamily\small, align=left, anchor=west] at (3.6,0.8) {\textcolor{FpgaBlue}{$\bullet$} positive: top bit 0};
  \node[font=\sffamily\small, align=left, anchor=west] at (3.6,0.1) {\textcolor{FpgaOrange}{$\bullet$} negative: top bit 1};
  \node[font=\sffamily\small, align=left, anchor=west] at (3.6,-0.8) {$7 + 1 = -8$: \textbf{overflow}!};
\end{tikzpicture}
\caption{Two's complement numbers are conveniently shown on a circle: adding
one is a step clockwise}
\end{figure}

The main advantage of two's complement: \textbf{one and the same adder can both
add and subtract}. To get $-b$, invert all the bits of $b$ and add one:
\[
  a - b = a + \overline{b} + 1.
\]
The inversion is done by XOR gates (with a control signal of 1, the gate
$x\oplus 1 = \overline{x}$ inverts), and the one is supplied by the carry input
of the lowest bit.

\begin{figure}[H]
\centering
\begin{tikzpicture}
  \node[hblock, minimum width=40mm, minimum height=14mm] (add) at (0,0) {4-bit\\adder};
  \node[block, minimum width=22mm, minimum height=10mm] (xr) at (1.0,2.2) {4 XOR\\\scriptsize gates};
  \draw[bus, -{Stealth}] (1.0,3.6) node[above, text=FpgaDark] {$B$} -- (xr.north);
  \draw[bus, -{Stealth}] (xr.south) -- node[right, lbl] {$B$ or $\overline{B}$} (1.0,0.7);
  \draw[bus, -{Stealth}] (-1.0,3.6) node[above, text=FpgaDark] {$A$} -- (-1.0,0.7);
  \draw[arr, FpgaRed] (-3.2,2.2) node[left, text=FpgaDark] {SUB} -- (xr.west);
  \draw[arr, FpgaRed] (-2.6,2.2) |- node[pos=0.75, above, lbl, text=FpgaRed] {$c_0$} (add.west);
  \fill[FpgaRed] (-2.6,2.2) circle (1.5pt);
  \draw[bus, -{Stealth}] (add.south) -- ++(0,-0.8) node[below, text=FpgaDark] {$S = A \pm B$};
\end{tikzpicture}
\caption{An adder-subtractor: with SUB $= 0$ it computes $A+B$, with SUB $= 1$
it computes $A-B$}
\end{figure}

\begin{diybox}[Subtraction]
\textbf{Task.} \stars{1} Compute $7-3$ by hand as $0111+1101$ and confirm in
Logra. \stars{2} Build an adder-subtractor as in the figure: four XOR gates on
the input $B$ and the carry $c_0=\text{SUB}$. \stars{3} Apply $3-5$ and explain
the negative result.

\textbf{In Logra.} Learn tab, the ``Arithmetic'' chapter, the lessons
``Subtraction'' and ``Changing the sign''; the reference is the example
\code{full-subtractor}.

\textbf{Expected result.} $7-3=0100=4$; $3-5=1110=-2$ (top bit 1).

\textbf{How to check.} Waveforms: with SUB $=0$ the output is $A+B$, with
SUB $=1$ the output is $A-B$; compare with the table of four-bit numbers in
two's complement.
\end{diybox}

\section{Overflow and Sign}

Four bits hold numbers from $-8$ to $+7$. If a result does not fit in this
range, the adder still outputs four bits, but they are \emph{wrong}. This is
\term{overflow}. For example:
\begin{gather*}
  5 + 3:\ 0101_2 + 0011_2 = 1000_2 = -8, \\
  (-8) + (-1):\ 1000_2 + 1111_2 = 1\,0111_2 \to 0111_2 = +7.
\end{gather*}
Two positive numbers gave a negative one, two negative numbers gave a positive
one. This is the sign of overflow in signed numbers: \emph{the addends have one
sign, and the result has the other}. Addends of different signs cannot overflow:
their sum lies between them.

The circuit detects overflow from the carries of the top bit. Let $c_3$ be the
carry \emph{into} the top bit and $c_4$ the carry \emph{out of} it. Then
\[
  V = c_3 \oplus c_4 .
\]
Let us check: in $5+3$ we have $c_3=1$, $c_4=0$, $V=1$; in $7-3$ ($0111+1101$)
$c_3=1$, $c_4=1$, $V=0$, and the carry $c_4=1$ is simply discarded---the result
$0100=4$ is correct. We see that \textbf{a carry out of the top bit and overflow
are different things}: the carry $c_4$ means going out of range only for
\emph{unsigned} numbers ($15+1$), while overflow $V$ does so for signed ones.

Next to the result, an ALU usually outputs four \term{flags}:
\begin{itemize}
  \item $Z$ (zero) --- all bits of the result are zero;
  \item $N$ (negative) --- the top bit of the result;
  \item $C$ (carry) --- $c_4$;
  \item $V$ (overflow) --- $c_3\oplus c_4$.
\end{itemize}

The flags give a comparison from the result of the subtraction $A-B$. For signed
numbers, $A<B$ if $N\oplus V=1$ (the sign of the result corrected for overflow).
For unsigned numbers, $A<B$ if the carry $C=0$: when subtraction is done as
addition of $\overline{B}+1$, the absence of a carry means a ``borrow''. Check
with $3-5$: $0011+1011=1110$, $c_4=0$, $N=1$, $V=0$---both tests say ``less''.

\begin{table}[H]
\centering
\caption{Flags in addition of four-bit numbers (signed values in parentheses)}
\small
\begin{tabular}{l c c c c c c}
\toprule
\textbf{Example} & \textbf{Result} & $Z$ & $N$ & $C$ & $V$ & \textbf{Correct?} \\
\midrule
$5+3$ (8) & \code{1000} & 0 & 1 & 0 & 1 & no, overflow \\
$7+(-3)$ (4) & \code{0100} & 0 & 0 & 1 & 0 & yes \\
$(-8)+(-1)$ (7) & \code{0111} & 0 & 0 & 1 & 1 & no, overflow \\
$3+(-3)$ (0) & \code{0000} & 1 & 0 & 1 & 0 & yes \\
\bottomrule
\end{tabular}
\end{table}

What to do on overflow is up to the designer: raise an error signal, take the
result with \emph{saturation} (show the maximum $7$ instead of $8$), or widen
the number. Signed numbers are widened by \term{sign extension}: the top bit is
copied into the new upper bits (for example, $1101_2=-3$ becomes
$1111\,1101_2=-3$).

\begin{diybox}[Overflow and sign]
\textbf{Task.} \stars{1} Add $0110_2+0101_2$ and explain why the result is
wrong. \stars{2} Build an adder with an output $V$ (\code{V = c3 xor c4}) and
check the table above. \stars{3} Make a saturating adder: when $V=1$ the result
is $0111$ (if the addends are positive) or $1000$ (if they are negative).

\textbf{In Logra.} Designer tab; the references are the examples
\code{carry-overflow} and \code{saturating}; to bring out the carries $c_3,c_4$
you need the chain of adders from the previous task.

\textbf{Expected result.} $6+5$ gives \code{1011} with $V=1$; $7+(-3)$ gives
\code{0100} with $V=0$.

\textbf{How to check.} Waveforms: go through the four examples of the flags
table and check $V$ and $C$.
\end{diybox}

\section{Comparators}

A \term{comparator} compares two numbers. The easiest thing to check is
\emph{equality}: the numbers are equal if all their bits are equal. Equality of
two bits is checked by an XNOR gate ($\overline{a\oplus b} = 1$ when $a = b$),
and an AND gate combines the results of all the bits.

\begin{figure}[H]
\centering
\begin{tikzpicture}
  \foreach \i in {0,...,3} {
    \node[gXNOR, minimum width=10mm, minimum height=7mm] (x\i) at (1.5,1.8-\i*1.2) {};
    \draw[wire] (x\i.input 1) -- ++(-0.6,0) node[left] {$a_\i$};
    \draw[wire] (x\i.input 2) -- ++(-0.6,0) node[left] {$b_\i$};
  }
  \node[gAND, logic gate inputs=nnnn, minimum height=18mm] (g) at (4.2,0) {};
  \foreach \i [evaluate=\i as \j using int(\i+1)] in {0,...,3}
    \draw[wire] (x\i.output) -- ++(0.3,0) |- (g.input \j);
  \draw[wire] (g.output) -- ++(0.6,0) node[right] {$A = B$};
\end{tikzpicture}
\caption{An equality comparator for 4-bit numbers}
\label{fig:cmp}
\end{figure}

Comparing ``greater/less'' is more involved: the numbers are compared starting
from the top bit. The first bit in which they differ determines the result. For
one bit:
\[
  (a > b) = a\,\overline{b}, \qquad (a < b) = \overline{a}\,b, \qquad (a = b) = \overline{a\oplus b}.
\]
By the way, $A < B$ can also be checked with a subtractor: if $A - B$ is
negative (the top bit of the result is 1, with no overflow), then $A < B$. The
exact rule with the flags $N$ and $V$ is given in the section ``Overflow and
Sign''.

\begin{diybox}[Security Panel: a threshold alarm]
\textbf{Task.} \stars{1} Build an equality comparator for 4-bit numbers as in
Fig.~\ref{fig:cmp}. \stars{2} Build an ``$A<B$'' block and check it on the pairs
$(3,5)$, $(5,3)$, $(5,5)$. \stars{3} Connect it up: $A$ is the noise level
(4 bits, set by switches), $B$ is the threshold; the alarm LED is on if the
level is \emph{greater} than the threshold.

\textbf{In Logra.} Learn tab, the ``Comparison and codes'' chapter, the lessons
``Equality in one block'' and ``Less than: A < B''; the ``Projects'' chapter, the
lesson ``A three-way comparator''; the references are the examples
\code{comparator} and \code{range-check}.

\textbf{Expected result.} The alarm is on if the level $A>B$ and off for
$A\le B$; when $A=B$ there is no alarm.

\textbf{Security Panel project milestone: milestone~5.} A threshold alarm.
Acceptance criterion: with the threshold $B=8$ the LED is on for levels
$9\ldots 15$ and off for $0\ldots 8$ (Waveforms tab, going through all 16
levels); changing the threshold moves the boundary.

\textbf{How to check.} Waveforms: of the outputs $A>B$, $A=B$, $A<B$, exactly one
is 1 for any inputs.
\end{diybox}

\section{Multiplication and the ALU}

Binary multiplication is even simpler than decimal: each digit of the
multiplier is 0 or 1, so the partial products are either zero or a shifted
multiplicand. All that remains is to add them.

\begin{figure}[H]
\centering
\begin{tikzpicture}[x=8mm, y=6.5mm, font=\ttfamily]
  \foreach \d [count=\i] in {1,0,1,1} \node at (\i+3,4) {\d};
  \foreach \d [count=\i] in {0,1,1,0} \node at (\i+3,3) {\d};
  \node at (2.8,3.5) {$\times$};
  \draw (3.3,2.5) -- (7.6,2.5);
  \foreach \d [count=\i] in {0,0,0,0} \node[text=FpgaGray] at (\i+3,2) {\d};
  \foreach \d [count=\i] in {1,0,1,1} \node[text=FpgaBlue] at (\i+2,1) {\d};
  \foreach \d [count=\i] in {1,0,1,1} \node[text=FpgaBlue] at (\i+1,0) {\d};
  \foreach \d [count=\i] in {0,0,0,0} \node[text=FpgaGray] at (\i,-1) {\d};
  \draw (0.5,-1.5) -- (7.6,-1.5);
  \foreach \d [count=\i] in {1,0,0,0,0,1,0} \node at (\i+0.0,-2) {\d};
  \node[font=\sffamily\small, anchor=west] at (8.5,4) {$1011_2 = 11$};
  \node[font=\sffamily\small, anchor=west] at (8.5,3) {$0110_2 = 6$};
  \node[font=\sffamily\small, anchor=west] at (8.5,-2) {$1000010_2 = 66$};
  \node[font=\sffamily\small, anchor=west, text=FpgaBlue] at (8.5,0.5) {shifted copies};
\end{tikzpicture}
\caption{Long multiplication: shifts and additions}
\end{figure}

An $n\times n$-bit multiplier requires about $n^2$ full adders, so it turns out
large. \lookahead{an FPGA has ready-made hardware multipliers for this (DSP
blocks, from digital signal processing; Chapter~\ref{ch:inside}).}

An adder, a subtractor, logic operations, and shifts, put together and selected
by a multiplexer according to an operation code, form an \term{arithmetic logic
unit} (ALU)---the heart of every processor.

\begin{figure}[H]
\centering
\begin{tikzpicture}
  \draw[thick, fill=FpgaLight] (0,1.8) -- (0.9,1.8) -- (1.35,1.1) -- (1.8,1.8) -- (2.7,1.8) -- (2.0,-0.4) -- (0.7,-0.4) -- cycle;
  \node[font=\sffamily\bfseries] at (1.35,0.5) {ALU};
  \draw[bus, -{Stealth}] (0.45,2.6) node[above, text=FpgaDark] {$A$} -- (0.45,1.8);
  \draw[bus, -{Stealth}] (2.25,2.6) node[above, text=FpgaDark] {$B$} -- (2.25,1.8);
  \draw[bus, -{Stealth}] (1.35,-0.4) -- (1.35,-1.2) node[below, text=FpgaDark] {result};
  \draw[arr, FpgaBlue] (-1.3,0.7) node[left, text=FpgaDark, align=right] {operation\\code} -- (0.33,0.7);
  \draw[arr] (2.37,0.7) -- (3.4,0.7) node[right, font=\sffamily\small, align=left] {flags: zero,\\carry, sign,\\overflow};
  \node[font=\sffamily\small, align=left, anchor=west] at (7.2,0.7) {\code{000}: $A+B$\\\code{001}: $A-B$\\\code{010}: $A \,\&\, B$\\\code{011}: $A \mid B$\\\code{100}: $A \oplus B$\ldots};
\end{tikzpicture}
\caption{The arithmetic logic unit}
\end{figure}

\begin{ideabox}
Addition is performed by full adders connected by a carry chain. Negative
numbers are written in two's complement, and then subtraction reduces to
addition. Overflow of a signed result is determined as $V=c_3\oplus c_4$. An
equality comparator is built from XNOR and AND gates, and ``greater/less'' is
built with a subtractor and flags.
\end{ideabox}

\begin{diybox}[A two-bit ALU]
\textbf{Task.} \stars{1} Go through the lesson and build a two-bit ALU with four
operations. \stars{2} Check each operation on several pairs $a,b$ and write down
a table of results. \stars{3} Add a fifth operation (for example, $A\oplus B$)
and an output $Z$ (the result is zero).

\textbf{In Logra.} Learn tab, the ``Projects'' chapter (the ``Circuits and the
board'' book), the lesson ``A two-bit ALU''; the reference is the example
\code{alu}.

\textbf{Expected result.} For each operation code $op$, the result of the
corresponding unit (adder, logic gates) passes to the output---this is the
multiplexer from Chapter~\ref{ch:mux}; the other units compute too, but their
results are not used.

\textbf{How to check.} Waveforms: with $a=2$, $b=1$ go through $op$ and compare
the outputs with the table; RTL: find the multiplexer that selects the result.
\end{diybox}

\section{Common Mistakes}

\begin{warnbox}[Common mistakes]
\begin{itemize}
  \item \textbf{Confusing carry and overflow.} With signed numbers the carry $c_4=1$ can occur with a correct result ($7+(-3)$); the sign of an error is $V=c_3\oplus c_4$. And conversely: $15+1$ in unsigned numbers gives an error, yet $V=0$.
  \item \textbf{Forgetting the ``$+1$'' in subtraction.} $A-B=A+\overline{B}+1$. If you apply only $A+\overline{B}$, the result will be one too small (for example, $6-2$ gives 3). The one is supplied by the carry $c_0$ of the lowest bit.
  \item \textbf{Comparing signed numbers as unsigned.} $1111_2$ as unsigned is 15 and is greater than $0001_2$, but as signed it is $-1$ and is smaller. The type of the numbers must be known before choosing a comparator.
  \item \textbf{Extending a signed number with zeros.} To turn $1101_2=-3$ into eight bits you must append ones ($1111\,1101$), not zeros ($0000\,1101=13$).
\end{itemize}
\end{warnbox}

\section{Questions and Problems}
\label{sec:tasks-arith}

Answers are collected in Appendix~\ref{app:answers}.

\begin{taskbox}[Chapter problems]
\begin{enumerate}
  \item \stars{1}~Add $1011_2 + 0110_2$ in binary and check the result in decimal.
  \item \stars{1}~Write the number $-5$ in 8-bit two's complement.
  \item \stars{1}~What is the range of numbers held by an 8-bit signed register in two's complement?
  \item \stars{1}~A full adder received $a=1$, $b=1$, $c_{in}=1$. What are $sum$ and $c_{out}$?
  \item \stars{1}~The carry delay of one bit is 0.5~ns. What is the delay of an 8-bit ripple-carry adder?
  \item \stars{2}~Compute $7 - 3$ in 4-bit two's complement, replacing subtraction with addition.
  \item \stars{2}~A comparator received $A=1010_2$ and $B=1001_2$. Which of the outputs ($A>B$, $A=B$, $A<B$) becomes one, and how can you see it from the bits?
  \item \stars{2}~Add the signed 4-bit numbers $0110_2+0101_2$. What is the result, what are $c_3$, $c_4$, $V$, and is the answer correct?
  \item \stars{2}~For $A=1101_2$ and $B=0010_2$ (signed), compute $A-B$ as $A+\overline{B}+1$ and use the flags $N$, $V$ to determine whether $A<B$ is true.
  \item \stars{2}~\emph{Find the error.} A student computes $6-2$ in 4 bits like this: $0110+\overline{0010}=0110+1101=1\,0011$, takes the low four bits and gets $0011=3$. Where is the mistake?
  \item \stars{2}~A four-bit ALU adds for code \code{000} and subtracts for code \code{001}. For $A=0111_2$, $B=0001_2$ and code \code{000} (the numbers are signed), find the result and the flags $Z,N,C,V$.
  \item \stars{3}~\textbf{Project: Security Panel milestone 5.} A threshold alarm: a 4-bit noise level $A$ and a threshold $B$; the lamp is on when $A>B$. \emph{Acceptance criterion:} with $B=8$ the lamp is on exactly for $A=9\ldots 15$ (going through 16 values on the Waveforms tab); with $B=15$ it never turns on.
  \item \stars{3}~\textbf{Project: ``Saturating adder''.} A 4-bit unsigned adder whose result is clamped to $1111$ when the carry $c_4=1$. \emph{Acceptance criterion:} $7+3=10$, $9+9=15$, $15+1=15$ (Waveforms tab; the reference is the example \code{saturating}).
  \item \stars{3}~\textbf{Project: ``ALU with flags''.} Extend the two-bit ALU: operations $A+B$, $A-B$, $A\,\&\,B$, $A\mid B$ and a flag $Z$ (the result is zero). \emph{Acceptance criterion:} for the pairs $(2,1)$, $(3,3)$, $(0,0)$ all four operations give the values from a hand-made table, and $Z=1$ only for the result $0$.
\end{enumerate}
\end{taskbox}

\section{What's Next}

All the circuits of this chapter and the previous ones are \emph{combinational}:
the output at each moment is determined only by the inputs at that moment.
Change an input and the output changes; remove it and the effect disappears.
They remember nothing: an adder does not know what it added a second ago, and a
comparator will not notice that an alarm happened and passed. But a security
panel must \emph{remember} an alarm after the door has been closed, a counter
must keep its value, and a calculator needs to store an intermediate result. So
far we have had circuits without memory; the next chapter (Chapter~\ref{ch:ff})
adds the key missing piece---an element that stores one bit.
